\documentclass[11pt,a4paper]{article}
\usepackage[T1]{fontenc}
\usepackage[utf8]{inputenc}
\usepackage[english]{babel}
\usepackage{amsmath,amssymb}
\usepackage[margin=2cm]{geometry}
\setlength{\parindent}{0pt}
\setlength{\emergencystretch}{3em}
\allowdisplaybreaks

\title{Tensors given by their components: true statements (45)}
\author{}
\date{}
\begin{document}
\maketitle
Statements whose hypotheses give a tensor by its components, with a table for two indices and a list for one, and contract it with the declared metric. All of them are true and must be proved, with a certificate that the checker accepts.

\begin{enumerate}
% Lists of numbers: scalar products, lowering an index (IBS winter camp, exercise 1).
\item $\displaystyle g_{ij} = \operatorname{diag}(1, 1) \land A^{i} = (1, 2) \land B_{i} = (3, 4) \implies A^{i} B_{i} = 11$
\item $\displaystyle g_{ij} = \operatorname{diag}(1, 1) \land A^{i} = (1, 2) \implies A^{i} A_{i} = 5$
\item $\displaystyle g_{ij} = \operatorname{diag}(1, 1) \land A^{i} = (1, 2) \land C_{i} = (1, 2) \implies \forall i \in \{1, 2\}\ \left( A_{i} = C_{i} \right)$
\item $\displaystyle g_{ij} = \operatorname{diag}(-1, 1) \land A^{i} = (1, 2) \implies A^{i} A_{i} = 3$
\item $\displaystyle g_{ij} = \operatorname{diag}(2, 3) \land A^{i} = (1, 2) \land C_{i} = (2, 6) \implies \forall i \in \{1, 2\}\ \left( A_{i} = C_{i} \right)$
\item $\displaystyle g^{ij} = \operatorname{diag}(\frac{1}{2}, \frac{1}{3}) \land A^{i} = (1, 2) \implies A^{i} A_{i} = 14$
% Without a metric: Cartesian components, with the Levi-Civita symbol or the Kronecker delta.
\item $\displaystyle a_{i} = (1, 2, 3) \land b_{i} = (4, 5, 6) \land c_{i} = (-3, 6, -3) \implies \varepsilon_{ijk} a_{j} b_{k} = c_{i}$
\item $\displaystyle a_{i} = (1, 0, 0) \land b_{i} = (0, 1, 0) \land c_{i} = (0, 0, 1) \implies \varepsilon_{ijk} a_{i} b_{j} c_{k} = 1$
\item $\displaystyle M_{ij} = \begin{pmatrix} 1 & 2 & 3 \\ 4 & 5 & 6 \\ 7 & 8 & 10 \end{pmatrix} \implies \varepsilon_{ijk} M_{1i} M_{2j} M_{3k} = -3$
\item $\displaystyle M_{ij} = \begin{pmatrix} 1 & 2 & 3 \\ 4 & 5 & 6 \\ 7 & 8 & 10 \end{pmatrix} \implies \delta_{ij} M_{ij} = 16$
\item $\displaystyle A_{i} = (a, 2 b, a b - \frac{1}{2}) \implies \varepsilon_{ijk} A_{j} A_{k} = 0$
\item $\displaystyle A_{i} = (a, 2 b, a b - \frac{1}{2}) \implies \delta_{ij} A_{i} A_{j} = a a + 4 b b + a a b b - a b + \frac{1}{4}$
\item $\displaystyle F^{ij} = \begin{pmatrix} 1 & 2 \\ 3 & 4 \end{pmatrix} \implies \delta_{ij} F^{ij} = F^{11} + F^{22}$
\item $\displaystyle F^{ij} = \begin{pmatrix} 1 & 2 \\ 3 & 4 \end{pmatrix} \implies \delta_{11} F^{12} = 2$
\item $\displaystyle R_{ij} = \begin{pmatrix} 0 & 1 & 0 \\ -1 & 0 & 0 \\ 0 & 0 & 1 \end{pmatrix} \land I_{ij} = \begin{pmatrix} a & d & 0 \\ d & b & 0 \\ 0 & 0 & c \end{pmatrix} \land J_{ij} = \begin{pmatrix} b & -d & 0 \\ -d & a & 0 \\ 0 & 0 & c \end{pmatrix} \implies \forall i, j \in \{1, 2, 3\}\ \left( R_{ik} R_{jl} I_{kl} = \delta_{ik} J_{kj} \right)$
% Tables of numbers with a metric: the trace in four ways, mixed components.
\item $\displaystyle g_{ij} = \operatorname{diag}(2, 3) \land T_{ij} = \begin{pmatrix} 1 & 2 \\ 3 & 4 \end{pmatrix} \implies T^{i}{}_{i} = \frac{11}{6}$
\item $\displaystyle g_{ij} = \operatorname{diag}(2, 3) \land T_{ij} = \begin{pmatrix} 1 & 2 \\ 3 & 4 \end{pmatrix} \implies T_{i}{}^{i} = \frac{11}{6}$
\item $\displaystyle g_{ij} = \operatorname{diag}(2, 3) \land T_{ij} = \begin{pmatrix} 1 & 2 \\ 3 & 4 \end{pmatrix} \implies g^{ij} T_{ij} = \frac{11}{6}$
\item $\displaystyle g_{ij} = \operatorname{diag}(2, 3) \land T_{ij} = \begin{pmatrix} 1 & 2 \\ 3 & 4 \end{pmatrix} \implies g_{ij} T^{ij} = \frac{11}{6}$
\item $\displaystyle g_{ij} = \operatorname{diag}(2, 3) \land T_{ij} = \begin{pmatrix} 1 & 2 \\ 3 & 4 \end{pmatrix} \land U^{i}{}_{j} = \begin{pmatrix} \frac{1}{2} & 1 \\ 1 & \frac{4}{3} \end{pmatrix} \implies \forall i, j \in \{1, 2\}\ \left( T^{i}{}_{j} = U^{i}{}_{j} \right)$
\item $\displaystyle g_{ij} = \operatorname{diag}(2, 3) \land T_{ij} = \begin{pmatrix} 1 & 2 \\ 3 & 4 \end{pmatrix} \implies T^{12} = \frac{1}{3}$
\item $\displaystyle g_{ij} = \operatorname{diag}(-1, 1, 4, 1) \land T_{ij} = \begin{pmatrix} 0 & 1 & 2 & 0 \\ 0 & s & 0 & 0 \\ c & 0 & 0 & 0 \\ 0 & 0 & 0 & 1 \end{pmatrix} \implies g^{ij} T_{ij} = 1 + s$
\item $\displaystyle g_{ij} = \operatorname{diag}(-1, 1, 4, 1) \land T_{ij} = \begin{pmatrix} 0 & 1 & 2 & 0 \\ 0 & s & 0 & 0 \\ c & 0 & 0 & 0 \\ 0 & 0 & 0 & 1 \end{pmatrix} \implies T^{i}{}_{i} = 1 + s$
\item $\displaystyle g_{ij} = \operatorname{diag}(-1, 1, 4, 1) \land T_{ij} = \begin{pmatrix} 0 & 1 & 2 & 0 \\ 0 & s & 0 & 0 \\ c & 0 & 0 & 0 \\ 0 & 0 & 0 & 1 \end{pmatrix} \implies T_{i}{}^{i} = 1 + s$
\item $\displaystyle g_{ij} = \operatorname{diag}(-1, 1, 4, 1) \land T_{ij} = \begin{pmatrix} 0 & 1 & 2 & 0 \\ 0 & s & 0 & 0 \\ c & 0 & 0 & 0 \\ 0 & 0 & 0 & 1 \end{pmatrix} \implies g_{ij} T^{ij} = 1 + s$
\item $\displaystyle g_{ij} = \operatorname{diag}(-1, 1, 4, 1) \land T_{ij} = \begin{pmatrix} 0 & 1 & 2 & 0 \\ 0 & s & 0 & 0 \\ c & 0 & 0 & 0 \\ 0 & 0 & 0 & 1 \end{pmatrix} \implies T^{3}{}_{1} = \frac{1}{4} c$
% The electromagnetic field tensor, with Greek indices from 0 to 3 (Carati, Baird, UMD, Dinius).
\item $\displaystyle \mu, \nu \in \{0, 1, 2, 3\} \land g_{\mu\nu} = \operatorname{diag}(1, -1, -1, -1) \land F_{\mu\nu} = \begin{pmatrix} 0 & E_{1} & E_{2} & E_{3} \\ -E_{1} & 0 & -H_{3} & H_{2} \\ -E_{2} & H_{3} & 0 & -H_{1} \\ -E_{3} & -H_{2} & H_{1} & 0 \end{pmatrix} \land G^{\mu\nu} = \begin{pmatrix} 0 & -E_{1} & -E_{2} & -E_{3} \\ E_{1} & 0 & -H_{3} & H_{2} \\ E_{2} & H_{3} & 0 & -H_{1} \\ E_{3} & -H_{2} & H_{1} & 0 \end{pmatrix} \implies \forall \mu, \nu \in \{0, 1, 2, 3\}\ \left( F^{\mu\nu} = G^{\mu\nu} \right)$
\item $\displaystyle \mu, \nu \in \{0, 1, 2, 3\} \land g_{\mu\nu} = \operatorname{diag}(1, -1, -1, -1) \land F_{\mu\nu} = \begin{pmatrix} 0 & E_{1} & E_{2} & E_{3} \\ -E_{1} & 0 & -B_{3} & B_{2} \\ -E_{2} & B_{3} & 0 & -B_{1} \\ -E_{3} & -B_{2} & B_{1} & 0 \end{pmatrix} \implies F_{\mu\nu} F^{\mu\nu} = 2 (B_{1} B_{1} + B_{2} B_{2} + B_{3} B_{3} - E_{1} E_{1} - E_{2} E_{2} - E_{3} E_{3})$
\item $\displaystyle \mu, \nu \in \{0, 1, 2, 3\} \land g_{\mu\nu} = \operatorname{diag}(1, -1, -1, -1) \land F_{\mu\nu} = \begin{pmatrix} 0 & E_{1} & E_{2} & E_{3} \\ -E_{1} & 0 & -B_{3} & B_{2} \\ -E_{2} & B_{3} & 0 & -B_{1} \\ -E_{3} & -B_{2} & B_{1} & 0 \end{pmatrix} \implies F^{01} = -E_{1}$
\item $\displaystyle \mu, \nu \in \{0, 1, 2, 3\} \land g_{\mu\nu} = \operatorname{diag}(1, -1, -1, -1) \land F_{\mu\nu} = \begin{pmatrix} 0 & E_{1} & E_{2} & E_{3} \\ -E_{1} & 0 & -B_{3} & B_{2} \\ -E_{2} & B_{3} & 0 & -B_{1} \\ -E_{3} & -B_{2} & B_{1} & 0 \end{pmatrix} \implies F^{12} = -B_{3}$
\item $\displaystyle \mu, \nu \in \{0, 1, 2, 3\} \land g_{\mu\nu} = \operatorname{diag}(1, -1, -1, -1) \land F^{\mu\nu} = \begin{pmatrix} 0 & -E_{1} & -E_{2} & -E_{3} \\ E_{1} & 0 & -B_{3} & B_{2} \\ E_{2} & B_{3} & 0 & -B_{1} \\ E_{3} & -B_{2} & B_{1} & 0 \end{pmatrix} \implies F_{01} = E_{1}$
\item $\displaystyle \mu, \nu \in \{0, 1, 2, 3\} \land g_{\mu\nu} = \operatorname{diag}(1, -1, -1, -1) \land F^{\mu\nu} = \begin{pmatrix} 0 & -E_{1} & -E_{2} & -E_{3} \\ E_{1} & 0 & -B_{3} & B_{2} \\ E_{2} & B_{3} & 0 & -B_{1} \\ E_{3} & -B_{2} & B_{1} & 0 \end{pmatrix} \implies F^{\mu\nu} F_{\mu\nu} = 2 (B_{1} B_{1} + B_{2} B_{2} + B_{3} B_{3}) - 2 (E_{1} E_{1} + E_{2} E_{2} + E_{3} E_{3})$
\item $\displaystyle \mu, \nu \in \{0, 1, 2, 3\} \land g_{\mu\nu} = \operatorname{diag}(1, -1, -1, -1) \land F^{\mu\nu} = \begin{pmatrix} 0 & -E_{1} & -E_{2} & -E_{3} \\ E_{1} & 0 & -B_{3} & B_{2} \\ E_{2} & B_{3} & 0 & -B_{1} \\ E_{3} & -B_{2} & B_{1} & 0 \end{pmatrix} \land \varepsilon^{0123} = 1 \land D^{\mu\nu} = \begin{pmatrix} 0 & -B_{1} & -B_{2} & -B_{3} \\ B_{1} & 0 & E_{3} & -E_{2} \\ B_{2} & -E_{3} & 0 & E_{1} \\ B_{3} & E_{2} & -E_{1} & 0 \end{pmatrix} \implies \forall \mu, \nu \in \{0, 1, 2, 3\}\ \left( D^{\mu\nu} = \frac{1}{2} \varepsilon^{\mu\nu\rho\sigma} F_{\rho\sigma} \right)$
\item $\displaystyle \mu, \nu \in \{0, 1, 2, 3\} \land g_{\mu\nu} = \operatorname{diag}(1, -1, -1, -1) \land F^{\mu\nu} = \begin{pmatrix} 0 & -E_{1} & -E_{2} & -E_{3} \\ E_{1} & 0 & -B_{3} & B_{2} \\ E_{2} & B_{3} & 0 & -B_{1} \\ E_{3} & -B_{2} & B_{1} & 0 \end{pmatrix} \land \varepsilon^{0123} = 1 \land D^{\mu\nu} = \begin{pmatrix} 0 & -B_{1} & -B_{2} & -B_{3} \\ B_{1} & 0 & E_{3} & -E_{2} \\ B_{2} & -E_{3} & 0 & E_{1} \\ B_{3} & E_{2} & -E_{1} & 0 \end{pmatrix} \implies D^{\mu\nu} D_{\mu\nu} = 2 (E_{1} E_{1} + E_{2} E_{2} + E_{3} E_{3}) - 2 (B_{1} B_{1} + B_{2} B_{2} + B_{3} B_{3})$
\item $\displaystyle \mu, \nu \in \{0, 1, 2, 3\} \land g_{\mu\nu} = \operatorname{diag}(1, -1, -1, -1) \land F^{\mu\nu} = \begin{pmatrix} 0 & -E_{1} & -E_{2} & -E_{3} \\ E_{1} & 0 & -B_{3} & B_{2} \\ E_{2} & B_{3} & 0 & -B_{1} \\ E_{3} & -B_{2} & B_{1} & 0 \end{pmatrix} \land \varepsilon^{0123} = 1 \land D^{\mu\nu} = \begin{pmatrix} 0 & -B_{1} & -B_{2} & -B_{3} \\ B_{1} & 0 & E_{3} & -E_{2} \\ B_{2} & -E_{3} & 0 & E_{1} \\ B_{3} & E_{2} & -E_{1} & 0 \end{pmatrix} \implies F^{\mu\nu} D_{\mu\nu} = -4 (E_{1} B_{1} + E_{2} B_{2} + E_{3} B_{3})$
\item $\displaystyle \mu, \nu \in \{0, 1, 2, 3\} \land g_{\mu\nu} = \operatorname{diag}(1, -1, -1, -1) \land \varepsilon^{0123} = 1 \land F_{\mu\nu} = \begin{pmatrix} 0 & E_{1} & E_{2} & E_{3} \\ -E_{1} & 0 & -B_{3} & B_{2} \\ -E_{2} & B_{3} & 0 & -B_{1} \\ -E_{3} & -B_{2} & B_{1} & 0 \end{pmatrix} \implies \varepsilon^{\mu\nu\lambda\rho} F_{\mu\nu} F_{\lambda\rho} = -8 (E_{1} B_{1} + E_{2} B_{2} + E_{3} B_{3})$
\item $\displaystyle \mu, \nu \in \{0, 1, 2, 3\} \land g_{\mu\nu} = \operatorname{diag}(1, -1, -1, -1) \land F^{\mu\nu} = \begin{pmatrix} 0 & -E_{1} & -E_{2} & -E_{3} \\ E_{1} & 0 & -B_{3} & B_{2} \\ E_{2} & B_{3} & 0 & -B_{1} \\ E_{3} & -B_{2} & B_{1} & 0 \end{pmatrix} \land L^{\mu}{}_{\nu} = \begin{pmatrix} c & s & 0 & 0 \\ s & c & 0 & 0 \\ 0 & 0 & 1 & 0 \\ 0 & 0 & 0 & 1 \end{pmatrix} \implies L^{0}{}_{\alpha} L^{2}{}_{\beta} F^{\alpha\beta} = -c E_{2} - s B_{3}$
\item $\displaystyle \mu, \nu \in \{0, 1, 2, 3\} \land g_{\mu\nu} = \operatorname{diag}(1, -1, -1, -1) \land F^{\mu\nu} = \begin{pmatrix} 0 & -E_{1} & -E_{2} & -E_{3} \\ E_{1} & 0 & -B_{3} & B_{2} \\ E_{2} & B_{3} & 0 & -B_{1} \\ E_{3} & -B_{2} & B_{1} & 0 \end{pmatrix} \land L^{\mu}{}_{\nu} = \begin{pmatrix} c & s & 0 & 0 \\ s & c & 0 & 0 \\ 0 & 0 & 1 & 0 \\ 0 & 0 & 0 & 1 \end{pmatrix} \implies L^{1}{}_{\alpha} L^{2}{}_{\beta} F^{\alpha\beta} = -c B_{3} - s E_{2}$
\item $\displaystyle \mu, \nu \in \{0, 1, 2, 3\} \land g_{\mu\nu} = \operatorname{diag}(1, -1, -1, -1) \land F^{\mu\nu} = \begin{pmatrix} 0 & -E_{1} & -E_{2} & -E_{3} \\ E_{1} & 0 & -B_{3} & B_{2} \\ E_{2} & B_{3} & 0 & -B_{1} \\ E_{3} & -B_{2} & B_{1} & 0 \end{pmatrix} \land L^{\mu}{}_{\nu} = \begin{pmatrix} c & s & 0 & 0 \\ s & c & 0 & 0 \\ 0 & 0 & 1 & 0 \\ 0 & 0 & 0 & 1 \end{pmatrix} \implies L^{0}{}_{\alpha} L^{1}{}_{\beta} F^{\alpha\beta} = -(c c - s s) E_{1}$
% The moment of inertia about a line (Vuk, Brescia), and a perfect fluid (MPE B.2).
\item $\displaystyle g_{ij} = \operatorname{diag}(1, 1, 1) \land I_{ij} = \begin{pmatrix} \frac{19}{3} & 0 & 0 \\ 0 & \frac{19}{6} & -\frac{485}{204} \\ 0 & -\frac{485}{204} & \frac{19}{6} \end{pmatrix} \land m^{i} = (0, 1, 1) \implies m^{i} I_{ij} m^{j} = \frac{161}{204} m^{k} m_{k}$
\item $\displaystyle g_{ij} = \operatorname{diag}(1, 1, 1) \land I_{ij} = \begin{pmatrix} 8 & -\frac{8}{3} & 0 \\ -\frac{8}{3} & \frac{7}{3} & 0 \\ 0 & 0 & \frac{31}{3} \end{pmatrix} \land m^{i} = (1, 1, 0) \implies m^{i} I_{ij} m^{j} = \frac{5}{2} m^{k} m_{k}$
\item $\displaystyle \mu, \nu \in \{0, 1, 2, 3\} \land g^{\mu\nu} = \operatorname{diag}(-1, 1, 1, 1) \land u^{\mu} = (c, 0, 0, 0) \land T^{\mu\nu} = \begin{pmatrix} r c c & 0 & 0 & 0 \\ 0 & P & 0 & 0 \\ 0 & 0 & P & 0 \\ 0 & 0 & 0 & P \end{pmatrix} \implies \forall \mu, \nu \in \{0, 1, 2, 3\}\ \left( c c T^{\mu\nu} = (r c c + P) u^{\mu} u^{\nu} + c c P g^{\mu\nu} \right)$
% Indices that start from 0, with Latin letters.
\item $\displaystyle i \in \{0, 1, 2\} \land A_{i} = (a, 2 b, 7) \implies \delta_{ij} A_{i} A_{j} = A_{0} A_{0} + A_{1} A_{1} + A_{2} A_{2}$
\item $\displaystyle i \in \{0, 1, 2\} \land A_{i} = (a, 2 b, 7) \implies \delta_{00} A_{0} = a$
\item $\displaystyle i \in \{0, 1, 2\} \land A_{i} = (a, 2 b, 7) \implies \delta_{00} A_{2} = 7$
\end{enumerate}
\end{document}
